\documentclass[12pt,a4paper]{article}
\usepackage[utf8]{inputenc}
\usepackage{amsmath,amssymb,amsthm}
\usepackage{algorithm}
\usepackage{algpseudocode}
\usepackage{geometry}
\geometry{margin=1in}

\newtheorem{theorem}{Theorem}
\newtheorem{lemma}[theorem]{Lemma}

\title{Problem 73: Counting Fractions in a Range}
\author{Project Euler Solution}
\date{}

\begin{document}
\maketitle

\section{Problem Statement}

How many reduced proper fractions lie strictly between $1/3$ and $1/2$ in the Farey sequence $F_{12000}$?

\section{Mathematical Analysis}

\begin{theorem}[Range Bounds]
For fixed $d \geq 2$, the integers $n$ with $1/3 < n/d < 1/2$ satisfy $\lfloor d/3 \rfloor + 1 \leq n \leq \lfloor (d-1)/2 \rfloor$.
\end{theorem}

\begin{proof}
$n/d > 1/3 \Rightarrow n > d/3 \Rightarrow n \geq \lfloor d/3 \rfloor + 1$. Similarly, $n/d < 1/2 \Rightarrow n < d/2 \Rightarrow n \leq \lfloor (d-1)/2 \rfloor$.
\end{proof}

\begin{theorem}[M\"obius Identity for Coprimality]
For positive integers $n, d$:
\[
[\gcd(n,d) = 1] = \sum_{k \mid \gcd(n,d)} \mu(k).
\]
\end{theorem}

\begin{proof}
Define $F(g) = \sum_{k \mid g} \mu(k)$. Since $\mu * \mathbf{1} = \varepsilon$ (the Dirichlet identity), $F(g) = \varepsilon(g) = [g = 1]$. Explicitly, for $g = p_1^{a_1} \cdots p_r^{a_r}$ with $g > 1$:
\[
F(g) = \prod_{i=1}^{r}\bigl(1 + (-1)\bigr) = 0,
\]
and $F(1) = 1$.
\end{proof}

\begin{theorem}[M\"obius Summation]
The count of reduced fractions in $(1/3, 1/2)$ with denominator $\leq N$ is:
\[
C(N) = \sum_{k=1}^{N} \mu(k) \cdot G\!\left(\left\lfloor N/k \right\rfloor\right),
\]
where $G(M) = \sum_{d=2}^{M} \bigl(\lfloor (d-1)/2 \rfloor - \lfloor d/3 \rfloor\bigr)$.
\end{theorem}

\begin{proof}
Expand the coprimality condition via Theorem~2:
\[
C(N) = \sum_{d=2}^{N} \sum_{\lfloor d/3 \rfloor < n < \lceil d/2 \rceil} \sum_{k \mid \gcd(n,d)} \mu(k).
\]
Substitute $d = kd'$, $n = kn'$ and exchange summation order. For fixed $k$, the inner sum counts integers $n'$ with $d'/3 < n' < d'/2$, which equals $\lfloor (d'-1)/2 \rfloor - \lfloor d'/3 \rfloor$ when non-negative. Summing over $d' = 2, \ldots, \lfloor N/k \rfloor$ gives $G(\lfloor N/k \rfloor)$.
\end{proof}

\section{Algorithm}

\begin{algorithm}[H]
\caption{Count fractions in $(1/3, 1/2)$ via M\"obius summation}
\begin{algorithmic}[1]
\Require $N = 12000$
\State Sieve $\mu[1..N]$
\State $\text{count} \gets 0$
\For{$k = 1$ \textbf{to} $N$}
    \If{$\mu[k] = 0$} \textbf{continue} \EndIf
    \State $M \gets \lfloor N/k \rfloor$
    \State $c \gets 0$
    \For{$d = 2$ \textbf{to} $M$}
        \State $\text{lo} \gets \lfloor d/3 \rfloor + 1$,\quad $\text{hi} \gets \lfloor (d-1)/2 \rfloor$
        \If{$\text{hi} \geq \text{lo}$}
            \State $c \gets c + \text{hi} - \text{lo} + 1$
        \EndIf
    \EndFor
    \State $\text{count} \gets \text{count} + \mu[k] \cdot c$
\EndFor
\State \Return count
\end{algorithmic}
\end{algorithm}

\section{Complexity}

\textbf{M\"obius sieve:} $O(N \log \log N)$.

\textbf{Double sum:} $\sum_{k=1}^{N} \lfloor N/k \rfloor = O(N \log N)$ iterations (harmonic sum).

\textbf{Time:} $O(N \log N)$. \quad \textbf{Space:} $O(N)$.

\section{Answer}

\[
\boxed{7295372}
\]

\end{document}
